\documentclass[12pt, a4paper]{article}

\usepackage[utf8]{inputenc}
\usepackage[T1]{fontenc}
\usepackage[brazil]{babel}
\usepackage{geometry}
\usepackage{amsfonts, amsmath, amssymb}
\usepackage{graphicx}
\usepackage{hyperref}
\usepackage{titlesec}
\usepackage{indentfirst}
\usepackage{float}
\usepackage{listings}
\usepackage{xcolor}

\geometry{a4paper, margin={3cm,2cm,2cm,3cm}}
\setlength{\parindent}{1.25cm}

% --- Configuração de Cores e Estilo para MATLAB ---
\definecolor{codegreen}{rgb}{0,0.6,0}
\definecolor{codegray}{rgb}{0.5,0.5,0.5}
\definecolor{codepurple}{rgb}{0.58,0,0.82}
\definecolor{backcolour}{rgb}{0.95,0.95,0.92}

\lstdefinestyle{mystyle}{
	language=Matlab,                % Define a linguagem
	backgroundcolor=\color{white},  % Cor de fundo (white ou backcolour)
	commentstyle=\color{codegreen}, % Cor dos comentários
	keywordstyle=\color{blue},      % Cor das palavras-chave (if, else, end)
	numberstyle=\tiny\color{codegray}, % Estilo dos números de linha
	stringstyle=\color{codepurple}, % Cor das strings
	basicstyle=\ttfamily\footnotesize, % Fonte tipo "máquina de escrever"
	breakatwhitespace=false,         
	breaklines=true,                 % Quebra automática de linha
	captionpos=b,                    % Posição da legenda (bottom)
	keepspaces=true,                 
	numbers=left,                    % Posição dos números (esquerda)
	numbersep=5pt,                  
	showspaces=false,                
	showstringspaces=false,
	showtabs=false,                  
	tabsize=2,
	frame=single,                    % Moldura ao redor do código
	% --- Suporte para acentos no código (UTF-8) ---
	literate={á}{{\'a}}1 {ã}{{\~a}}1 {é}{{\'e}}1 {ç}{{\c{c}}}1 {ê}{{\^e}}1 {ó}{{\'o}}1 {í}{{\'i}}1 {ú}{{\'u}}1 {à}{{\`a}}1
}

\lstset{style=mystyle}

% --- Configuração de Títulos (com titlesec) ---
% Força os títulos a terem 12pt (\normalsize) e negrito

\titleformat{\section}
  {\normalsize\bfseries\MakeUppercase} % formato: 12pt e negrito
  {\thesection.}         % rótulo (número): "1."
  {1em}                  % separação: 1 "M" de espaço
  {}                     % código antes do título (pode deixar vazio)

\titleformat{\subsection}
  {\normalsize\MakeUppercase} % formato: 12pt e negrito
  {\thesubsection.}      % rótulo (número): "1.1."
  {1em}                  % separação
  {}                     % código antes

\titleformat{\subsubsection}
  {\normalsize\bfseries} % formato: 12pt e negrito
  {\thesubsubsection.}   % rótulo (número): "1.1.1."
  {1em}                  % separação
  {}                     % código antes

\newcommand{\universidade}{Universidade Federal de Santa Catarina}
\newcommand{\centro}{Centro Tecnológico de Joinville - CTJ}
\newcommand{\disciplina}{EMB5602 - Controle Digital}
\newcommand{\curso}{Engenharia Mecatrônica}
\newcommand{\professor}{Prof. Alexandro Garro Brito}
\newcommand{\titulo}{Controle de profundidade de veículo subaquático autônomo (AUV)}
\newcommand{\autor}{Hebert Alan Kubis}
\newcommand{\matricula}{23102547}

\newcommand{\vehicleMass}{M_t}
\newcommand{\vehicleVolume}{V_t}
\newcommand{\SensorPressH}{h_p}

\begin{document}
	\begin{titlepage}
		\begin{center}
			{\bfseries \MakeUppercase{\universidade}}\\
			{\bfseries \MakeUppercase{\curso}}\\[3cm]
			
			{\MakeUppercase{\autor}}\\

			\vfill
			{\bfseries \MakeUppercase{\titulo}}\\
			\vfill

			{Joinville}\\
			{\number\year}

		\end{center}
	\end{titlepage}
	
	% Folha de rosto (ABNT)
	\begin{titlepage}
		\begin{center}
			{\MakeUppercase{\autor}}\\
			\vfill

			{\bfseries \MakeUppercase{\titulo}}\\[2cm]
			
			\hspace*{8cm}
			\parbox[t]{\dimexpr\textwidth-8cm\relax}{
				{Trabalho apresentado ao \centro\ da\ \universidade\ como requisito 
				parcial para a disciplina \disciplina.}\\[1cm]
				{\professor}\\
			}\\

			\vfill

			{Joinville}\\
			{\number\year}

		\end{center}
	\end{titlepage}

	\begin{abstract}
		Fazer
		
		\vspace{1em}
		\noindent\textbf{Palavras-chave:} AUV, Controle de Profundidade, Estabilidade
	\end{abstract}
	\clearpage

	\tableofcontents
	\clearpage

	\section{Introdução}

	Fazer uma introdução

	\section{Definição do sistema}

	O sistema sob análise consiste em um Veículo Autônomo Subaquático (AUV -- \textit{Autonomous Underwater Vehicle}) impulsionado por um conjunto de seis propulsores. A configuração de atuadores permite a movimentação no plano horizontal ($xOy$) por meio de quatro thrusters U2 mini e movimento no eixo vertical ($Z$) proporcionado por dois thrusters T100 da BlueRobotics.

	Estruturalmente, o veículo é caracterizado por dois compartimentos cilíndricos estanques, alinhados longitudinalmente ao eixo $X$. O cilindro superior, de maior volume, aloja os sistemas embarcados e a eletrônica de controle. O cilindro inferior, de menor diâmetro, é destinado ao armazenamento da bateria, contribuindo para a estabilidade passiva do sistema (centro de gravidade abaixo do centro 
	de empuxo).
	
	Diante desta configuração, o objetivo central deste projeto de controle é definido como: 
	\textbf{garantir a estabilização do veículo em uma profundidade de referência especificada.}

		\subsection{Hipóteses de Modelagem}
			A modelagem matemática do sistema foi obtida com base em algumas hipóteses simplificadoras, que visam capturar a dinâmica essencial do controle de profundidade, ao mesmo tempo em que mantêm a complexidade do modelo em um nível gerenciável para análise e projeto de controladores. As hipóteses adotadas são as seguintes:
			
			\begin{itemize}
				\item \textbf{Desacoplamento dos eixos:} A dinâmica vertical foi analisada isoladamente (1 Grau de Liberdade), assumindo-se que o movimento nos planos horizontal e lateral não influencia significativamente a dinâmica de profundidade;

				\item \textbf{Atuação ideal:} Assumiu-se que a força resultante dos dois propulsores verticais atua diretamente sobre o centro de massa do veículo, não gerando torques de arfagem (\textit{pitch}) ou rolagem (\textit{roll}) indesejados;

				\item \textbf{Sistema de Coordenadas:} Adotou-se o referencial com o eixo $Z$ positivo apontando para baixo (padrão NED - \textit{North-East-Down}), onde o aumento de $Z$ corresponde ao aumento da profundidade;

				\item \textbf{Hidrostática constante:} O veículo é considerado completamente submerso durante toda a operação, garantindo que o volume deslocado e, consequentemente, a força de empuxo, permaneçam constantes.
			\end{itemize}

		\subsection{Instrumentação e Sensoriamento}
			A estratégia de controle em malha fechada depende da aquisição precisa dos estados do sistema. O AUV dispõe dos seguintes sensores e métodos para a realimentação:
			
			\begin{itemize}
				\item \textbf{Sensor de Pressão:} Responsável pela medição da pressão hidrostática, a partir da qual a profundidade linear ($z$) é estimada;

				\item \textbf{Unidade de Medida Inercial (IMU):} Integrada à controladora Pixhawk PX4, fornece dados de aceleração linear ($\ddot{z}$) e velocidade angular;

				\item \textbf{Estimação de Velocidade:} A velocidade vertical ($\dot{z}$) não é medida diretamente, sendo obtida através de técnicas de fusão sensorial (como Filtro de Kalman ou Complementar) que combinam a integração da aceleração com a derivação da posição.
			\end{itemize}

			Esses são os dados disponíveis para utilização no projeto, apesar de que o dado fundamental aqui é a profundidade $z$ medida pela Pixhawk atravez de uma combinação dos dados do sensor de pressão e da IMU, utilizando um filtro de Kalman para obter uma estimativa robusta e precisa da profundidade atual do veículo.

			
			FIZ ATE AQUI

	\section{Modelagem Matemática do Sistema}

		\subsection{Variáveis do Sistema}

			Para o desenvolvimento do modelo matemático de controle de profundidade do AUV, adotam-se variáveis 
			genéricas para representar os parâmetros físicos do veículo. Esta abordagem assegura a robustez e a 
			adaptabilidade do modelo frente a eventuais alterações no projeto mecânico. Os parâmetros fundamentais 
			são:

			\begin{itemize}
				\item $\vehicleMass$: Massa total do veículo;
				
				\item $\vehicleVolume$: Volume total deslocado pelo veículo;
			\end{itemize}

		\subsection{Análise Dinâmica e Equilíbrio de Forças}

			A modelagem inicia-se pela análise das forças verticais atuantes sobre o corpo do veículo submerso:

			\begin{itemize}
				\item Força Peso ($F_{p}$);
				
				\item Força de Empuxo ($F_{e}$);
				
				\item Força de Atuação dos Propulsores ($F_{m}$);
				
				\item Força de Arrasto Hidrodinâmico ($F_{d}$).
			\end{itemize}

			Aplicando a Segunda Lei de Newton para o movimento de translação no eixo vertical ($Z$), obtém-se o 
			somatório de forças:
			
			\begin{equation}
				\sum F = F_{p} - F_{e} + F_{m} - F_{d} = m \cdot a
			\end{equation}

			Onde $a = \frac{d^2z}{dt^2}$ representa a aceleração vertical do veículo. Substituindo as expressões 
			correspondentes a cada força:

			\begin{equation}
				M_t \cdot \frac{d^2z}{dt^2} = M_t \cdot g - \rho \cdot g \cdot V_t + F_{m} - B_d \cdot \frac{dz}{dt}
				\label{eq:dinamica_nao_linear}
			\end{equation}

			Onde $g$ é a aceleração da gravidade, $\rho$ é a densidade do fluido e $B_d$ é o coeficiente de 
			amortecimento viscoso linearizado.

\section{Conclusão}

	O presente trabalho abordou o desenvolvimento completo do sistema de controle de profundidade para o AUV Yvy, 
	abrangendo desde a modelagem física até a validação do controlador em frequência. A obtenção do modelo 
	matemático linearizado, fundamentado na física do arrasto hidrodinâmico e nos parâmetros reais do veículo, 
	revelou uma planta de segunda ordem do Tipo 1, intrinsecamente capaz de anular erros de regime para entradas 
	em degrau.

	A definição rigorosa dos requisitos de desempenho, priorizando a segurança operacional (sobressinal limitado 
	a 5\%), guiou o projeto do controlador pelo método do Lugar das Raízes. A análise comparativa entre as 
	estratégias de controle demonstrou que, embora a ação derivativa (PD) permitisse tempos de resposta inferiores 
	a 1 segundo, tal agressividade implicaria em esforços de controle excessivos e sensibilidade a ruídos. 
	Portanto, concluiu-se que o controlador puramente \textbf{Proporcional (P)}, com ganho ajustado para 
	$K_p = 30,35$, representa a solução mais equilibrada, garantindo um tempo de acomodação satisfatório 
	($T_s \approx 3,5$ s) e robustez contra incertezas de modelagem.

	Por fim, a análise da resposta em frequência validou a estabilidade e a robustez do projeto. O sistema 
	apresentou Margem de Ganho infinita e uma Margem de Fase elevada ($\approx 73^\circ$), o que assegura um 
	comportamento bem amortecido e tolerância a atrasos de transporte superiores a 1,8 segundos. A largura de 
	banda obtida ($\approx 1$ rad/s) confirmou que o sistema de malha fechada atuará efetivamente no rastreamento 
	de trajetórias de profundidade, filtrando naturalmente perturbações de alta frequência e ruídos de sensores, 
	atendendo assim a todos os objetivos de engenharia propostos para a operação do AUV Yvy.

	\clearpage
	\bibliographystyle{plain}
	\bibliography{referencias}
	
	\clearpage
	\appendix

	\section{Scripts de simulação e análise}

\end{document}